\author{Janne Marike Lorenzen}
\title{Individual variability in the encoding and decoding of prosodic prominence relations}
\subtitle{}
\BackTitle{Individual variability in the encoding and decoding of prosodic prominence relations}


\renewcommand{\lsID}{611}
\renewcommand{\lsSeries}{silp}
\renewcommand{\lsSeriesNumber}{18}
\ISBNdigital{978-3-96110-590-8}
\ISBNhardcover{978-3-98554-207-9}
\BookDOI{10.5281/zenodo.21133333}

\BackBody{This book investigates inter-individual variability in the encoding and decoding of prosodic prominence relations. Formally, prosodic prominence relations are redundantly expressed by several acoustic parameters, such as changes in fundamental frequency, duration, and intensity. Functionally, prominence relations encode meaning, e.g., the information status of a word or referent. Both the formal and the functional side of prominence facilitate inter-individual variability. On the one hand, redundancy inherent in prominence marking enables individuals to use different parameters in their production and perception. On the other hand, differences in communicative skills may influence which meaning aspects language users pay attention to and how strongly they encode them.

In this book, inter-individual variability is analyzed in two production studies and one perception study. Production data is elicited in listener-directed reading tasks and analyzed using prosodic annotation and acoustic analyses. Perception is probed via a prominence rating task. Individual differences are explored via careful descriptive inspection of the data and statistically via cluster analyses. Results indicate that speakers differ in how strongly they encode prominence relations, which parameters they use to do so, and which meaning aspects they predominantly mark - which in turn is partially modulated by their communicative skills. Listeners also vary in their sensitivity to prominence cues.

These findings highlight the importance of looking at the individual speaker or listener in prosodic research and allow for a new perspective on established concepts in prosodic theory, such as
the dichotomy of tune and metrical structure, relationality and redundancy in
prominence encoding and decoding, and the probabilistic nature of form-meaning mapping.
}

\renewcommand{\lsImpressumExtra}{The work presented in this book is based on the author's doctoral dissertation, which was accepted by the Faculty of Arts and Humanities at the University of Cologne in 2024.}

\typesetter{Janne Marike Lorenzen}
\proofreader{Matthew Korte}

